\section{Conclusion}
\label{sec:conclusion}

Alert triage asks a local model to do two jobs, reading probe results and deciding what to do with them, and small models are good at only the first. \hesp{} divides the jobs: the model reads, and a controller chooses the probes, decides when to stop, and accepts only verdicts that ledger evidence supports. On three simulated settings with a given list of candidate causes, a controller without any model resolves as many cases as every model we tested, whereas a fixed parser fails completely once the log format changes and a 7B reader does not. Reading is also where attackers write, so \hesp{} never lets the model's own reading support a benign verdict; with this rule, the finish guard, and source-diverse corroboration, no injected instruction, forged source, or adaptive attack on the reader closes a case. Whether a model should also share the decisions when the candidate causes are not known in advance, or when evidence only adds up in combination, requires settings that a fixed procedure cannot solve and data that the public security datasets we examined do not provide.
